\par\medskip
\begin{center}\hypertarget{native-163-article116}{}\pdfbookmark[2]{116.}{article-116}116.\end{center}
Before undertaking to determine an orbit from geocentric observations, we must apply certain reductions to them for nutation, precession, parallax, and aberration if the utmost precision is required. These small corrections can be neglected in a rougher calculation.

Observations of planets and comets are usually given in apparent right ascensions and declinations, that is, referred to the apparent position of the equator. Since this position varies because of nutation and precession and therefore differs from one observation to another, it is advisable first to introduce a fixed plane in place of the variable one. For this purpose we may adopt either the equator in its mean position for some epoch or the ecliptic. The latter plane is usually used, but the former also has particular advantages that deserve attention.

Whenever the plane of the equator is chosen, the observations must first be corrected for nutation and then reduced to an arbitrary epoch by applying precession. This operation is exactly the same as the one by which the mean position of a fixed star for a given epoch is derived from its observed position, and therefore needs no explanation here. If we decide to adopt the plane of the ecliptic, two methods are available. We may use the mean obliquity to deduce the longitudes and latitudes from right ascensions and declinations corrected for nutation and precession, in which case the resulting longitudes will already be referred to the mean equinox. Or, more conveniently, we may use the apparent obliquity to compute the longitudes and latitudes from apparent right ascensions and declinations and then correct the longitudes for nutation and precession.

The positions of the Earth corresponding to the individual observations are computed from the solar tables. They must clearly be referred to the same plane as the observations of the celestial body. Accordingly, nutation is neglected in computing the longitude of the Sun. This longitude is then reduced by applying precession to the \SourceLeaf{164}fixed epoch and increased by 180 degrees. If it seems worthwhile to take the Sun's latitude into account, we give it the opposite sign. This gives the heliocentric position of the Earth. If the equator has been chosen as the fundamental plane, this position can be transformed into right ascension and declination by means of the mean obliquity.
